\documentclass[]{article}

\usepackage{amsmath}
\usepackage{multirow}
\usepackage{natbib}
\usepackage{graphicx} 
\usepackage[margin=1in]{geometry}

\begin{document}

The nature of anomalous diffusion in two-dimensional turbulence remains a subject of significant interest, particularly regarding how Lagrangian transport statistics evolve during the inverse energy cascade. In this study, we investigate the relationship between wavelet-filtered eddy lifetimes and Lévy flight characteristics by re-advecting tracers in velocity fields derived from 1024² direct numerical simulations. We analyze the temporal evolution of the Lévy stability exponent $\alpha$ and its stationarity as the energy spectrum undergoes condensation toward the largest scales. Our analysis reveals that $\alpha$ is non-stationary and scale-dependent, decreasing from approximately 2.0 at smaller scales to 1.75 at the largest scales. While we observe power-law tails in residence time and flight displacement distributions, our results indicate that the observed non-stationarity of $\alpha$ is not directly driven by spectral condensation or the spatial density of hyperbolic points. These findings suggest that the complexity of Lagrangian dynamics in two-dimensional turbulence arises from multi-scale interactions that are not fully captured by simple continuous-time random walk models or instantaneous spectral evolution.

\end{document}
                